\section{Results}

\subsection{A minimal retrieval-key and instructed-CA1 architecture}

The model contains four functional components (Fig.~\ref{fig:model}A).
Entorhinal input, denoted $x_{\mathrm{EC}}$, projects through a fixed sparse
map to a CA3-like population that provides a retrieval key. A separately
pretrained EC--CA1 encoder transforms the same input into an instructed CA1
state (IS). Plastic CA3--CA1 weights associate the retrieval key with the IS,
whereas the corresponding pretrained CA1--EC decoder remains fixed. During
recall, the EC input regenerates a key, the learned association reconstructs a
CA1 state, and the decoder maps that state back to EC output. This architecture
separates the selection of associative content from its expression in a
pre-existing output basis.

We first evaluated two plasticity rules under matched inputs and a shared
reference parameter configuration
(Fig.~\ref{fig:model}B). The instructive-driven rule writes an outer-product
association gated by the IS without considering the current recalled state.
The bounded error-driven rule instead potentiates CA3--CA1 synapses for underexpressed
CA1 components and depresses them for overexpressed components. Repeated
presentations therefore rewrite the instructed association under the instructive-driven
rule but correct a residual under the error-driven rule. Both are abstract
models of rapid instructed plasticity; neither implements the temporal
eligibility trace or dendritic plateau dynamics of biological BTSP.

Track inputs combined a MEC-like spatial component with a LEC-like cue
component (Fig.~\ref{fig:model}D). Two cue identities occurred at fixed
positions and could exchange locations across blocks. For degradation tests,
the model first stored intact laps; plasticity was then frozen while a fixed
fraction of either the MEC-like or LEC-like input units was set to zero
(Fig.~\ref{fig:model}E).

\begin{figure}[!htbp]
    \centering
    \includegraphics[width=\textwidth]{figure_1.png}
    \caption{\textbf{Model architecture and simulation protocol.}
    \textbf{(A)} EC input produces a fixed sparse CA3-like retrieval key
    $x_{\mathrm{CA3}}$ and, through a pretrained encoder, an instructed CA1
    state $x_{\mathrm{IS}}$. Plastic
    CA3--CA1 weights store their association; a fixed pretrained CA1--EC
    decoder reads out recall. \textbf{(B)} Instructive-driven and bounded
    error-driven update rules.
    \textbf{(C)} Conceptual reconstruction-error trajectories over first and
    repeated presentations; colors denote pattern identity. This panel is a
    schematic rather than a measured result. \textbf{(D)} Circular-track input
    composed of MEC-like spatial activity and LEC-like cue activity.
    \textbf{(E)} Selective degradation protocol. Intact laps are stored before
    plasticity is frozen; a fraction of MEC-like or LEC-like units is then
    zeroed using a mask held fixed within a probe lap and resampled across
    probes.}
    \label{fig:model}
\end{figure}

\subsection{Readout compatibility constrains rapid associative storage}

We first isolated the relationship between an instructed CA1 code and its fixed
output decoder (Fig.~\ref{fig:compatibility}A). Twenty independently
initialized autoencoders were pretrained on structured MEC+LEC patterns and
then frozen. Pattern order was shuffled, so this assay retained the marginal
spatial--cue structure without presenting a traversed track sequence.
For each initialization, 28 held-out patterns were stored once under five
paired conditions. In the aligned condition, the IS remained in the
coordinates learned by the decoder. In the fixed-permutation condition, the
same non-identity permutation was applied to every IS while the decoder was
left unchanged. In the matched-decoder condition, the corresponding decoder
columns were permuted together with the IS. Random-content and no-plasticity
conditions controlled for storing an incorrect association and for leaving the
memory unwritten, respectively.

Under the instructive-driven rule, aligned instructions produced a mean output--target
cosine of 0.880 (95\% confidence interval [CI], 0.870--0.891), whereas a fixed
coordinate permutation reduced it to 0.247 (0.207--0.286). The paired
aligned-minus-permuted difference was 0.634 (0.592--0.676). The same coordinate
dependence appeared under the error-driven rule: aligned and permuted means
were 0.872 (0.858--0.887) and 0.178 (0.147--0.209), respectively, giving a
paired difference of 0.694 (0.658--0.731). Permuting the decoder together with
the IS restored the aligned mean exactly for each rule. Random-content recall
remained lower (0.355 for the instructive-driven rule and 0.315 for the
error-driven rule), as did the shared
no-plasticity control (0.114). The matched-decoder restoration is an algebraic
coordinate control: the permutation preserves the IS values but changes their
meaning to the unchanged decoder. Reconstruction therefore depends on
compatibility between the instructed CA1 coordinates and the fixed output
basis, not on information content alone. Under this reference configuration,
the two rules reached similar aligned one-shot recall despite using different
update signals: the instructive-driven rule performs a full instructed write,
whereas the error-driven rule makes a bounded residual correction.

\subsection{Cue exchanges induce reproducible CA1 reconfiguration}

We next exposed the model to 40 track laps in which two cue identities exchanged
positions every 10 laps (Fig.~\ref{fig:compatibility}B). A paired no-swap
control used the same root seeds, pretrained encoder and decoder, MEC-like
trajectories, lap count, plasticity parameters, and exact EC--CA3 projection,
but retained the initial cue assignment throughout.
After every training lap, plasticity was paused and the model was probed in the
context scheduled for that lap.

For the instructive-driven rule, mean unit-wise tuning similarity between consecutive probes
fell to 0.365 at cue exchanges (95\% CI, 0.348--0.382), compared with 0.845 at
the same lap boundaries in the no-swap control (0.832--0.857). The paired
swap-minus-control difference was $-0.479$ ($-0.498$ to $-0.461$). Under the
error-driven rule, exchange-boundary similarity was 0.451 (0.436--0.466),
compared with 0.827 in the matched control (0.817--0.838), for a paired
difference of $-0.376$ ($-0.396$ to $-0.356$). At non-exchange transitions in
the alternating schedule, similarity remained high for both the instructive-driven
rule (0.833, 0.823--0.843) and the error-driven rule (0.821, 0.814--0.828).
The abrupt tuning changes were
therefore associated with the altered cue sequence rather than with continued
learning alone.

We next aligned the three exchange events to their first post-exchange lap
(Fig.~\ref{fig:compatibility}E). Under the instructive-driven rule,
consecutive-lap similarity at the exchange was 0.375 (95\% CI,
0.351--0.398), versus 0.851 (0.839--0.863) at matched no-swap boundaries; the
paired difference was $-0.477$ ($-0.499$ to $-0.454$). Under the error-driven
rule, the corresponding values were 0.462 (0.443--0.480) and 0.833
(0.822--0.845), with a paired difference of $-0.372$ ($-0.391$ to $-0.352$).
By the following lap, similarity recovered to 0.769 (0.756--0.783) and 0.819
(0.806--0.832), respectively. Thus the response was concentrated at the cue
exchange rather than developing as nonspecific drift over the subsequent
block.

The lap-by-lap population representational-similarity matrix provided a
complementary view (Fig.~\ref{fig:compatibility}F). Within-block maps remained
similar, exchanges separated the sequence into distinct blocks, and
nonadjacent blocks carrying the same cue assignment resembled one another.
The model therefore alternated between cue-dependent population states rather
than accumulating a wholly novel map after every block. Retrospective
comparison with the final map of each new block further showed that the first
post-exchange representation was already closer to the eventual adapted map
than to the pre-exchange map, followed by smaller within-block changes
(Supplementary Fig.~\ref{fig:remapping-dynamics}A,B). Endpoint similarities
were broadly distributed across CA1 units, demonstrating that the population
effect was not produced by a single uniform response class
(Supplementary Fig.~\ref{fig:remapping-dynamics}C,D).

After learning, plasticity was frozen and both cue arrangements were probed.
The representative error-driven population contained spatially stable fields,
amplitude-modulated responses, and fields whose preferred position changed
with cue identity (Fig.~\ref{fig:compatibility}C). Across seeds, mean
unit-wise, mean-centered spatial stability was 0.330 (95\% CI, 0.297--0.363)
under the instructive-driven rule and 0.464 (0.441--0.486) under the
error-driven rule. Mean absolute cue modulation
was 0.132 (0.130--0.135) and 0.124 (0.120--0.128), respectively. In both rules,
the unit-level distributions were continuous rather than separated into
discrete response classes (Fig.~\ref{fig:compatibility}D), and strong cue
modulation was concentrated among units with low or negative cross-context
stability. These distributions are descriptive: units nested within a seed
were not treated as independent replicates, and the analysis is not a formal
classification of rate versus global remapping. Operationally, changing the
cue arrangement changed both the EC-derived instruction and the CA3-like key,
producing a mixture of stable and context-dependent recalled CA1 activity.

\begin{figure}[!htbp]
    \centering
    \includegraphics[width=0.98\textwidth]{figure_2.pdf}
    \caption{\textbf{Readout compatibility and cue-dependent CA1
    reconfiguration.} \textbf{(A)} Output--target cosine for shuffled MEC+LEC
    samples in aligned,
    fixed-permutation, matched-decoder, random-content, and no-plasticity
    conditions for the instructive-driven rule (left) and error-driven rule
    (right). Grey points are
    seed means across 28 stored memories; black points and error bars show mean
    $\pm$ SEM across 20 seeds. \textbf{(B)} Mean unit-wise tuning similarity
    between consecutive scheduled-context probes for the instructive-driven
    rule (left) and error-driven rule (right). Purple denotes the cue-exchange
    schedule and grey the matched
    no-swap control; vertical dashed lines mark exchanges. Curves show mean
    $\pm$ SEM across 20 paired seeds. \textbf{(C)} Recalled CA1 activity on the
    same track under cue contexts A (left) and B (right) for the representative
    error-driven seed. Units have the same order and are sorted by preferred position
    across contexts; color limits are shared. \textbf{(D)} Unit-level
    relationship between mean-centered spatial stability and cue modulation
    for the instructive-driven rule (left) and error-driven rule (right),
    pooling 50 units from each of 20 seeds
    (1,000 points per rule). The vertical line marks zero stability. Panels C
    and D illustrate heterogeneous stable and cue-dependent responses but do
    not formally classify remapping type; confidence intervals in the text use
    seed-level summaries. \textbf{(E)} Event-aligned consecutive-lap tuning
    similarity for the instructive-driven rule (left) and error-driven rule
    (right) in the cue-exchange schedule and exact-wiring no-swap control. Lap
    zero is the first post-exchange lap;
    values were averaged across three exchanges within each seed and curves
    show mean $\pm$ SEM across 20 seeds. \textbf{(F)} Error-driven
    lap-by-lap population-map similarity, averaged across 20 seeds. Before
    comparison, each unit was mean-centered over track position and the full
    position-by-unit map was vectorized. Dashed lines mark cue-exchange
    boundaries. Panels B and E used identical EC--CA3 projections in the
    exchange and no-swap schedules for each rule; panel F shows the
    error-driven exchange schedule.}
    \label{fig:compatibility}
\end{figure}

\subsection{Selective EC degradation dissociates cue and spatial readout}

We then asked whether the stored association remained usable when one EC input
modality was selectively corrupted (Fig.~\ref{fig:degradation}). Each model
stored 20 intact laps with a fixed cue arrangement. Plasticity was frozen, and
independent masks removed 0, 25, 50, 75, 90, or 100\% of eligible units. Each
fraction used 12 masks per seed and 20 paired root seeds; the numerical summary
below focuses on 90\% partial degradation. The principal control
replaced the normal sparse EC--CA3 map with a dense common map in which every
CA3-like unit received the same average EC input. This manipulation removes
item-specific key structure and is therefore a degenerate-key, not a
sparsity-matched, control.

Under LEC-like degradation, cue identity accuracy for the normal key declined
from 1.000 to 0.871 at 90\% dropout under the instructive-driven rule (95\% CI
at 90\%, 0.803--0.938) and from 1.000 to 0.829 under the error-driven
rule (0.766--0.893;
Fig.~\ref{fig:degradation}A). The corresponding 90\%-dropout changes were
$-0.129$ ($-0.197$ to $-0.062$) and $-0.171$ ($-0.234$ to $-0.107$).
MEC output--target cosine remained near its instructive-driven baseline (0.736 at
baseline and 0.723 at 90\% dropout; change $-0.013$, $-0.029$ to
$0.003$), whereas it declined from 0.904 to 0.780 under the error-driven rule
(change $-0.125$,
$-0.136$ to $-0.113$; Fig.~\ref{fig:degradation}B). Thus spatial output
remained substantially readable under LEC-like corruption, although it was not
strictly invariant under the error-driven rule.

MEC-like degradation produced the complementary loss of position information.
Nearest-position accuracy declined from 0.163 to 0.070 under the
instructive-driven rule (90\%-dropout
change $-0.093$, 95\% CI $-0.106$ to $-0.079$) and from 0.666 to 0.069 under
the error-driven rule (change $-0.597$, $-0.636$ to $-0.557$; chance $=0.02$;
Fig.~\ref{fig:degradation}C). Cue identity remained 1.000 for every seed and
degradation level under both rules (Fig.~\ref{fig:degradation}D). The low
baseline position accuracy under the instructive-driven rule is itself
informative: repeated
error-driven correction yielded a substantially more discriminable spatial
readout before degradation, whereas the instructive-driven rule did not. The two rules
should therefore not be interpreted as starting from equivalent clean spatial
performance.

The dense common-key control remained at chance for cue identity (0.5) and
position (0.02) even without degradation, and produced low spatial-output
cosine. Its failure establishes that a non-specific common key cannot support
the associations, but does not isolate sparsity from key discriminability.
Together, the normal-key results demonstrate degraded-cue robustness and
pathway-selective failure within this architecture. They do not establish
recurrent pattern completion: the CA3-like layer has no attractor dynamics, and
recall deteriorates once the feedforward retrieval key loses discriminability.
Supplementary Fig.~\ref{fig:rule-comparison} collects four seed-paired contrasts
that summarize how the rule choice changes these effects.

\subsection{Evolution identifies rule-specific operating regimes and
load-dependent trade-offs}

The matched reference configuration above isolates consequences of changing
the update equation, but it need not be favorable to both rules. We therefore
performed separate covariance matrix adaptation evolution strategy (CMA-ES)
searches for the instructive-driven and error-driven rules on three fixed
development seeds. Within each generated MEC+LEC lap, storage samples were
shuffled before learning, preserving the input distribution while removing
adjacent track states. The objective equally weighted clean recall, recall after
85\% MEC-like degradation, and recall after 85\% LEC-like degradation; these
development evaluations were not treated as experimental replicates. To keep
the representational basis common across the principal and optimization-informed
analyses, these searches fixed $K_{\mathrm{CA1}}=5$ and both autoencoder gains
at 25. The independently selected medial-temporal-lobe (MTL) model
configurations used fan-in 36,
$K_{\mathrm{CA3}}=6$, $\beta_{\mathrm{CA3}}=286.747$,
$\beta_{\mathrm{CA1}}=9.160$, and $\alpha=0.03366$ for the
instructive-driven rule, versus fan-in 9, $K_{\mathrm{CA3}}=5$,
$\beta_{\mathrm{CA3}}=54.400$, $\beta_{\mathrm{CA1}}=26.928$, and
$\alpha=0.08197$ for the error-driven rule. Their development-set composite
scores were 0.714 and 0.818, respectively. Because both parameter vectors were selected on the same
three seeds used to define the objective, these values describe optimization
outcomes rather than unbiased estimates of generalization.

A one-parameter-at-a-time analysis evaluated multipliers from 0.1 to 9 around
each selected vector while fixing all other parameters
(Fig.~\ref{fig:degradation}E). Ten independent evaluation seeds, disjoint from
the search seeds, supplied paired score differences and their across-seed
standard deviations. The composite score changed smoothly near most selected
values but was not uniformly insensitive. The largest mean decreases were
produced by increasing $K_{\mathrm{CA3}}$ ninefold for instructive-driven
plasticity ($-0.231$) and reducing $\beta_{\mathrm{CA3}}$ tenfold for
error-driven plasticity ($-0.477$). Small positive mean differences at a few
tested points were at most 0.002 and were small relative to their across-seed
variation. The analysis therefore supports a local
operating region while revealing important parameter asymmetries; it is not a
global robustness guarantee.

We next tested the selected configurations on 10 held-out seeds while
increasing associative load from 1 to all 90 non-repeating ordered pairs drawn
from 10 cue identities (Fig.~\ref{fig:degradation}F). Every pair was presented
three times, and exact two-cue recall was measured from clean probes and after
50\% LEC-unit dropout. Instructive-driven clean recall declined from
$1.000\pm0.000$ at one pair to $0.210\pm0.043$ at 10 pairs and
$0.138\pm0.016$ at 90 pairs (mean $\pm$ SEM); corresponding degraded recall
was $1.000\pm0.000$, $0.035\pm0.012$, and $0.051\pm0.006$. Error-driven clean
recall was $1.000\pm0.000$, $0.390\pm0.046$, and $0.151\pm0.022$ at the same
loads, with degraded values of $0.775\pm0.095$, $0.098\pm0.024$, and
$0.046\pm0.003$. Thus both selected regimes exhibited interference as load
increased, with an error-driven advantage at intermediate clean load but
convergence toward low exact-pair recall at the largest load. These curves are
empirical load profiles for this cue vocabulary, training schedule, and model
size rather than universal capacity estimates.

\begin{figure}[!htbp]
    \centering
    \includegraphics[width=\textwidth]{figure_3.pdf}
    \caption{\textbf{Selective degradation, parameter sensitivity, and
    cue-pair load.} In panels A--D, each pair of plots shows the
    instructive-driven rule on the left and the error-driven rule on the right.
    \textbf{(A)} Cue-identity accuracy during
    LEC-only degradation.
    \textbf{(B)} MEC output--target cosine during LEC-only degradation.
    \textbf{(C)} MEC nearest-position accuracy during MEC-only degradation.
    \textbf{(D)} Cue-identity accuracy during MEC-only degradation. Curves show
    mean $\pm$ SEM across 20 seeds after averaging 12 masks within each seed.
    Blue denotes the normal sparse CA3-like retrieval map; green denotes the
    dense common-key control, which collapses item-specific key structure and
    is not sparsity matched. Dotted horizontal lines denote chance accuracy
    where applicable. All degraded probes were performed with plasticity
    disabled. \textbf{(E)} Change in the three-component composite recall
    score when one MTL parameter was multiplied by the indicated factor while
    all others were held at that rule's evolved value. Integer parameters were
    rounded and model-valid bounds were enforced; each cell shows the paired
    mean score change over 10 independent seeds and its across-seed standard
    deviation. \textbf{(F)} Exact recall of
    ordered cue pairs as the number stored increased from 1 to 90. Blue denotes
    the instructive-driven rule and red the error-driven rule; open solid curves use intact probes and filled dashed
    curves use 50\% LEC dropout. Points show mean $\pm$ SEM across 10 held-out
    seeds, and the dotted line marks chance ($1/100$). Panels E and F use
    separately evolved rule-specific configurations and therefore compare
    operating regimes, whereas panels A--D use the shared reference
    configuration to isolate the update rule.}
    \label{fig:degradation}
\end{figure}

\FloatBarrier
